% Generated by scripts/build-paper.py — edit paper/src/survey.src.md
\section{Scope, definitions and what is documented}\label{sec:scope-definitions-and-what-is-documented}

\subsection{The interface}

A caller sends a \texttt{state} (a string, JSON object or array of text values) and a map of typed questions; the model evaluates every question independently against the shared state and returns one answer per question \citep{typesafe2026state,typesafe2026primitives}. Three primitives exist. A \textbf{Choice} selects among named options, each defined by a description; the answer includes the chosen option, a probability for every option and a \texttt{confidence} value. The launch post states support for up to 255 options, with a two-stage procedure for high cardinality \citep{typesafe2026launch}. A \textbf{Score} places the answer on two to ten ordered, described levels; it returns a position that may fall between levels, probabilities per level and \texttt{confidence}, and the vendor warns against interpolating exact magnitudes from it \citep{typesafe2026jaggedness}. A \textbf{Noul} returns the probability that a yes/no proposition holds and carries no separate confidence \citep{typesafe2026primitives,typesafe2026confidence}.

At the cutoff the documented model was still \texttt{jev-1.13.0}; the aliases \texttt{jev-latest} and \texttt{jev-preview} both pointed to it, and the vendor advises pinning the versioned ID when thresholds are tuned to a version \citep{typesafe2026models}. Input is text only; images, audio and video must first be converted to text or structured fields, and English is the primary training language \citep{typesafe2026models,typesafe2026state}. Visual and audio Jev-like models in the open ecosystem are independent projects, not TypeSafe releases.

\subsection{Probability semantics}

Probabilities and confidence are different objects. The documentation defines \texttt{confidence} as a statistic computed from the returned distribution and recommends thresholds tested on the user's own data \citep{typesafe2026confidence}; the exact statistic is not published. Empirically, one judging study found that native confidence detected errors about as well as the largest label probability (AUROC 0.745–0.876 across three benchmarks) \citep{arxiv2609_26550}, and a biosecurity audit read it as a chance-normalised maximum probability \citep{arxiv2609_30454}. The vendor states that calibration is measured across groups of predictions and does not guarantee that an individual answer is correct \citep{typesafe2026systemone}. Two independent sources observe that probabilities are returned on a 0.01 grid, which bounds calibration error for rare events from below and can put exactly zero on a correct answer \citep{arxiv2609_24052,gh_scienthoon_jev_ood_calibration}.

\subsection{What is documented as weak, and what is not disclosed}

The vendor maintains a list of known weak spots of jev-1.13. On its 2 October revision the list names literal reading, math and numbers (including counting and numeric representations), date and time comparison, indirection, large states full of irrelevant detail, adversarial content, contradictory instructions and criteria, Choice option order — the model “leans toward the option that comes first” — and generation \citep{typesafe2026jaggedness}. Several studies below measure exactly these edges. The architecture, training data and objective are not published; the documentation says one set of weights serves every account \citep{typesafe2026models}. Community reconstructions exist and remain speculation about the commercial model \citep{gh_jaredpalmer_kev}.

\subsection{Vendor claims as hypotheses}

The launch headline “193.6× faster, 444.6× cheaper” comes from the vendor's workflow evaluations, whose reference answers average two external frontier models and whose workflows were written by its capabilities team \citep{typesafe2026launch}. The vendor's public evaluation dashboard now averages four structured workflows against reference labels from two frontier models: in workflow mode it lists Jev at 67.8\% agreement for \$0.0004 and 0.4 s per case, against 74.1\% for the best listed configuration and 73.1\% for Opus 5, and every listed model agrees more often when the task is decomposed into a workflow than when it is asked as one prompt \citep{typesafe2026evals}. The plotted 0\% type-error rate is analytical rather than empirical \citep{typesafe2026launch}. We record these figures as vendor claims, not as evaluations.

\subsection{Names that collide}

Several names mean different things. TypeSafe's \emph{System One} borrows Kahneman's System 1 metaphor, which AI research had already adopted \citep{arxiv2010_06002,arxiv2502_17419}. \emph{RLCD} in the Jev context expands to Reinforcement Learning for Calibrated Decisions and is unrelated to Reinforcement Learning from Contrastive Distillation (2023), a distinction one core study also draws \citep{arxiv2609_24574}; open reimplementations of the Jev meaning now exist \citep{arxiv2609_38850,arxiv2610_02486}. \emph{OpenJev}, \emph{Open-Jev} and \emph{open-jev} name at least eight unrelated projects: frozen-decoder readouts, among them SemIf under its former name, two DiffusionGemma readouts, a family of fine-tuned decision models, an in-browser runtime, and the open model of one paper's title, JevLite \citep{gh_theoleecj_semif,gh_daseinlabs_open_jev,gh_zhihz_openjev,gh_razorback16_openjev,gh_joshuasp_open_jev,gh_zefan_cai_open_jev,gh_nico_martin_open_jev,arxiv2609_23959}. Adjacency in name implies neither the same algorithm nor copying.
